\scanpage{31}
\begin{center}-- 213 --\end{center}
\hfill n$^\circ$ 391

\uline{11.} \uline{Système de type} $E_7$ :

(I)\quad Gardons les notations du n$^\circ$ 10, et considérons l'intersection $R'$
de $R$ et de l'hyperplan $V'$ de $V$ orthogonal au vecteur
$\varepsilon_7+\varepsilon_8$ ; ce sont donc les vecteurs
$x=\sum_{i=1}^8\xi_i\varepsilon_i$ de $R$ tels que $\xi_7+\xi_8=0$,
autrement dit :
\[
\pm\varepsilon_i\pm\varepsilon_j,\quad
(i\neq j,\ 1\leq i\leq6,\ 1\leq j\leq6),\quad
\pm(\varepsilon_7-\varepsilon_8)\text{ et}
\]
\[
\pm\frac12(\varepsilon_7-\varepsilon_8+
\sum_{i=1}^6(-1)^{\unsure{\nu|\upsilon}(i)}\varepsilon_i),
\quad\text{où }\sum_{i=1}^6\unsure{\nu|\upsilon}(i)\text{ est \uline{impair}.}
\]
Il est clair que $R'$ engendre $V'$, donc est un système de racines réduit
de rang 7, ayant $n=2+\binom62.4+2^6=126$ éléments.

(II)\quad Munissons $V'$ de l'ordre induit par l'ordre lexicographique sur $V$ ;
les racines de $R'$ qui sont positives pour cet ordre sont, comme on l'a vu au
n$^\circ$ 10, de la forme $\sum_{i=1}^8n_i\alpha_i$, où les $n_i$ sont des entiers
$\geq0$ ; si l'on exprime qu'une telle racine est orthogonale à
$\varepsilon_7+\varepsilon_8$, il vient, (puisque la plus grande racine
$\varepsilon_7+\varepsilon_8$ de $R$ est orthogonale à tous les $\alpha_i$
sauf $\alpha_8$), $n_8=0$. Toutes les racines positives de $R'$,
(pour l'ordre considéré), sont donc combinaisons linéaires à coefficients positifs
des $\alpha_i$ pour $1\leq i\leq7$, ce qui montre que ces racines,
(qui sont orthogonales à $\varepsilon_7+\varepsilon_8$, donc appartiennent à $R'$),
forment un système de \uline{racines simples} $S'$ de $R'$.
Le calcul des $(\alpha_i|\alpha_j)$ fait au n$^\circ$ 10 montre alors que le
schéma de Dynkin de $R'$ est de type $E_7$, et $R'$ est irréductible.

(III)\quad En vertu du \S\ 8, n$^\circ$ 3, formule (4), on a $h=n/\ell=18$.

(IV)\quad La plus grande racine pour l'ordre induit par celui de $R^8$
défini au n$^\circ$ 10 est maintenant
$\varepsilon_8-\varepsilon_7=2\alpha_1+2\alpha_2+3\alpha_3+4\alpha_4+
3\alpha_5+2\alpha_6+\alpha_7$, donc c'est la plus grande racine
$\tilde\alpha$, (\S\ 6, n$^\circ$ 6, cor. 2 de la prop. 13),
et on vérifie que
\mspage{30}
